% Transcription of Cayley's Collected Mathematical Papers, Volume I, pages 76-100
% Paper 12 (conclusion): On the Theory of Determinants
% Paper 13: On the Theory of Linear Transformations
% Paper 14 (start): On Linear Transformations
%
% Method: pdftotext OCR base + scan PNG visual restoration at 200 dpi.
% Source scan: collectedmathema01cayluoft.pdf (Internet Archive). Public domain.

\documentclass[11pt,a4paper]{article}
\usepackage[utf8]{inputenc}
\usepackage[T1]{fontenc}
\usepackage{amsmath,amssymb,amsthm}
\usepackage{mathtools}
\usepackage{geometry}
\usepackage{fancyhdr}
\usepackage{titlesec}
\usepackage{array}
\usepackage{booktabs}
\geometry{margin=1in}
\theoremstyle{plain}
\newtheorem{theorem}{Theorem}
\newtheorem{lemma}{Lemma}
\theoremstyle{definition}
\newtheorem{definition}{Definition}
\theoremstyle{remark}
\newtheorem{remark}{Remark}
\pagestyle{fancy}
\fancyhf{}
\fancyhead[C]{\small CAYLEY---COLLECTED MATHEMATICAL PAPERS, VOL.~I}
\fancyfoot[C]{\thepage}
\renewcommand{\headrulewidth}{0pt}
\titleformat{\section}{\normalfont\Large\bfseries\centering}{}{0em}{}
\titleformat{\subsection}{\normalfont\large\bfseries}{\textsc{Chap.}~\thesubsection.}{0.5em}{}

\setcounter{section}{11}
\usepackage{graphicx}
\emergencystretch=3em
\tolerance=600
\begin{document}
\setlength{\parindent}{1.5em}
\setlength{\parskip}{0pt}
\thispagestyle{empty}
\newpage
\setcounter{page}{76}

% ============================================================
% PAPER 12 (conclusion): ON THE THEORY OF DETERMINANTS
% Pages 76--79
% ============================================================

\noindent\textit{[Paper 12 continued from p.~75.]}

\vspace{6pt}

This being premised, consider the symbol
\[
\left\{
\begin{array}{c}
A\rho\sigma\,\ldots(n) \\[2pt]
\hline
\rho_k\sigma_k\,\ldots
\end{array}
\right\}
\tag{4},
\]
denoting the sum of all the different terms of the form
\[
\pm_r\pm_s\,\ldots\,A\rho_{r_1}\sigma_{s_1}\,\ldots\,A\rho_k\sigma_k\,\ldots
\tag{5},
\]
the letters
\[
r_1,\ r_2\,\ldots r_k;\quad s_1,\ s_2\,\ldots s_k;\ \&\text{c.}
\tag{6},
\]
denoting any permutations whatever, the same or different, of the series of numbers
(2) [and the several combinations of $\rho$, $\sigma$ \ldots being understood as denoting suffixes of
the $A$'s]. The number of terms represented by the symbol (5) is evidently
\[
(1\cdot 2\,\ldots k)^n \tag{7}.
\]

In some cases it will be necessary to leave a certain number of the vertical
rows $\rho$, $\sigma$ \ldots unpermuted. This will be represented by writing the mark $(\dagger)$ immediately
above the rows in question. So that for instance
\[
\left\{
\begin{array}{c}
\stackrel{\dagger}{A\rho}\,\stackrel{\dagger}{\sigma}\,\ldots(n) \\[2pt]
\hline
\rho_k\sigma_k\,\ldots
\end{array}
\right\}
\tag{8},
\]
the number of rows with the $(\dagger)$ being $x$, denotes the sum of the
\[
(1\cdot 2\,\ldots k)^{n-x} \tag{9}
\]
terms, of the form
\[
\pm_r\pm_s\,\ldots\,A\rho_{r_1}\sigma_{s_1}\,\ldots\,A\rho_k\sigma_k\,\ldots
\tag{10}.
\]
Then it is obvious, that if all the rows have the mark $(\dagger)$ the notation (8) denotes
a single product only, and if the mark $(\dagger)$ be placed over all but one of the rows
the notation (8) belongs to a determinant. It is obvious also that we may write
\[
\left\{
\begin{array}{c}
\stackrel{\dagger}{A\rho}\,\stackrel{\dagger}{\sigma}\,\ldots(n) \\[2pt]
\hline
\rho_k\sigma_k\,\ldots
\end{array}
\right\}
= \Sigma\,\pm_u\pm_v\,\ldots
\left\{
\begin{array}{c}
A\rho_{u_1}\sigma_{v_1}\,\ldots\,A\rho_{u_k}\sigma_{v_k} \\[2pt]
\hline
\\
\end{array}
\right\}
\tag{11},
\]
where $\Sigma$ refers to the different permutations,
\[
u_1,\ u_2,\ldots u_k;\quad v_1,\ v_2,\ldots v_k;\ \&\text{c.}
\tag{12},
\]
which can be formed out of the numbers (2). The equation (11) would still be true,
if the mark $(\dagger)$ were placed over any number of the columns $\rho$, $\sigma$ \ldots

Suppose in this equation a single column only is left without the mark $(\dagger)$ on
the second side of the equation; the first side is then expressed as the sum of a number
\[
(1\cdot 2\,\ldots k)^{n-1},\quad\text{or generally}\quad (1\cdot 2\,\ldots k)^{n-x-1}\tag{13},
\]
of determinants, according as we consider the symbol (4) or the more general one (8).
And this may be done in $n$ or $(n-x)$ different ways respectively.

\newpage

It may be remarked, that the symbol (8) is the same in form as if a single
column only had the mark $(\dagger)$ over it; the number $n$ being at the same time
reduced from $n$ to $(n-x+1)$: for the marked columns of symbols may be replaced
by a single marked column of new symbols. Hence, without loss of generality, the
theorems which follow may be stated with reference to a single marked column only.

Suppose the letters
\[
\rho_1,\ \rho_2,\ldots \rho_k;\quad \sigma_1,\ \sigma_2,\ldots \sigma_k;\ \&\text{c.}
\tag{14}
\]
denote certain permutations of
\[
\alpha_1,\ \alpha_2,\ldots \alpha_k;\quad \beta_1,\ \beta_2,\ldots \beta_k;\ \&\text{c.},
\tag{15},
\]
in such a manner that
\[
\rho_1=\alpha_{g_1},\ \rho_2=\alpha_{g_2},\ldots\rho_k=\alpha_{g_k};\quad
\sigma_1=\beta_{h_1},\ \sigma_2=\beta_{h_2},\ldots\sigma_k=\beta_{h_k},\ \&\text{c.}
\tag{16}.
\]
Then the two following theorems may be proved:
\[
\left\{
\begin{array}{c}
\stackrel{\dagger}{A\rho}\,\sigma\,\ldots(n) \\[2pt]
\hline
\rho_k\sigma_k\,\ldots
\end{array}
\right\}
= \pm_g\pm_h\,\ldots
\left\{
\begin{array}{c}
\stackrel{\dagger}{A\alpha}\,\beta\,\ldots(n) \\[2pt]
\hline
\alpha_k\beta_k\,\ldots
\end{array}
\right\}
\tag{17},
\]
if $n$ be even; but in the contrary case
\[
\left\{
\begin{array}{c}
\stackrel{\dagger}{A\rho}\,\sigma\,\ldots(n) \\[2pt]
\hline
\rho_k\sigma_k\,\ldots
\end{array}
\right\}
= +\pm_g\,\ldots
\left\{
\begin{array}{c}
\stackrel{\dagger}{A\alpha}\,\beta\,\ldots(n) \\[2pt]
\hline
\alpha_k\beta_k\,\ldots
\end{array}
\right\}
\tag{18}.
\]

By means of these, and the equation (11), a fundamental property of the symbol
(8) may be demonstrated. We have
\[
\left\{
\begin{array}{c}
A\rho\sigma\,\ldots(n) \\[2pt]
\hline
\rho_k\sigma_k\,\ldots
\end{array}
\right\}
= \Sigma\,\pm_g\pm_h\,\ldots
\left\{
\begin{array}{c}
\stackrel{\dagger}{A\alpha}\,\beta\,\ldots(n) \\[2pt]
\hline
\alpha_k\beta_k\,\ldots
\end{array}
\right\}
\tag{19},
\]
which, when $n$ is even, reduces itself by (17) to
\[
\left\{
\begin{array}{c}
\stackrel{\dagger}{A\alpha}\,\stackrel{\dagger}{\beta}\,\ldots(n) \\[2pt]
\hline
\alpha_k\beta_k\,\ldots
\end{array}
\right\}
\Sigma(\pm_g\,1)
=1\cdot 2\,\ldots k
\left\{
\begin{array}{c}
\stackrel{\dagger}{A\alpha}\,\beta\,\ldots(n) \\[2pt]
\hline
\alpha_k\beta_k\,\ldots
\end{array}
\right\}
\tag{20}.
\]
But when $n$ is odd, from the equation (18),
\[
\left\{
\begin{array}{c}
\stackrel{\dagger}{A\alpha}\,\beta\,\ldots(n) \\[2pt]
\hline
\alpha_k\beta_k\,\ldots
\end{array}
\right\}
\Sigma(\pm_g\,1) = 0
\tag{21},
\]
since the number of negative and positive values of $\pm_g$ are equal.

\newpage

But when $n$ is odd, from the equation (18), since the number of negative and positive
values of $\pm_g$ are equal,
\[
\left\{
\begin{array}{c}
A\rho\sigma\,\ldots(n) \\[2pt]
\hline
\rho_k\sigma_k\,\ldots
\end{array}
\right\}
=
\left\{
\begin{array}{c}
\stackrel{\dagger}{A\alpha}\,\beta\,\ldots(n) \\[2pt]
\hline
\alpha_k\beta_k\,\ldots
\end{array}
\right\}
\Sigma(\pm_g\,1)=0
\tag{21}.
\]

From the equation (20), it follows that when $n$ is even, the value of a symbol
of the form
\[
\left\{
\begin{array}{c}
\stackrel{\dagger}{A\alpha}\,\beta\,\ldots(n) \\[2pt]
\hline
\alpha_k\beta_k\,\ldots
\end{array}
\right\}
\tag{22}
\]
is the same, over whichever of the columns $\alpha$, $\beta$ \ldots the mark $(\dagger)$ is placed. To
denote this indifference, the preceding quantity is better represented by
\[
\left\{
\begin{array}{c}
\stackrel{\dagger}{A\alpha}\,\stackrel{\dagger}{\beta}\,\ldots(n) \\[2pt]
\hline
\alpha_k\beta_k\,\ldots
\end{array}
\right\}
\tag{23},
\]
this last form being never employed when $n$ is odd, in which case the same property
does not hold. Hence also an ordinary determinant is represented by
\[
\left\{
\begin{array}{c}
\stackrel{\dagger}{A\alpha}\,\stackrel{\dagger}{\beta}\,\ldots \\[2pt]
\hline
\alpha_k\beta_k\,\ldots
\end{array}
\right\},
\quad
\left\{
\begin{array}{c}
\stackrel{\dagger}{A}\,1\,1\,\ldots \\[2pt]
\hline
k\,k
\end{array}
\right\}
\tag{24},
\]
the latter form being obviously equally general with the former one.

It is obvious from the equations (17), (18), that the expression (22) vanishes, in
the case of $n$ even whenever any two of the symbols $\alpha$ are equivalent, or any two
of the symbols $\beta$, \&c.; but if $n$ be odd, this property holds for the symbols $\beta$, \&c.,
but not for the marked ones $\alpha$. In fact, the interchange of the two equal symbols,
in each case, changes the sign of the expression (22), but they evidently leave it
unaltered, i.e.\ the quantity in question must be zero.

Consider now the symbol
\[
\left\{
\begin{array}{c}
\stackrel{\dagger}{A}\,1\,1\,\ldots(2p) \\[2pt]
\hline
k\,k
\end{array}
\right\}
\tag{25},
\]
which, for shortness, may be denoted by
\[
\{A\cdot k\cdot 2p\} \tag{26}.
\]

\newpage

I proceed to prove a theorem, which may be expressed as follows:
\[
\{A\cdot k\cdot 2p\}\cdot\{B\cdot k\cdot 2q\}=\{AB\cdot k\cdot 2p+2q-2\}
\tag{27},
\]
where
\[
\overline{AB}_{rs\ldots\,xy\ldots}=\Sigma\cdot A_{rs\ldots l}\,B_{xy\ldots l}
\tag{28},
\]
the number of the symbols $r$, $s$, \ldots being obviously $2p-1$, and that of $x$, $y$, \ldots being
$2q-1$. The summatory sign $\Sigma$ refers to $l$, and denotes the sum of the several terms
corresponding to values of $l$ from $l=1$ to $l=k$. Also the theorem would be equally
true if $l$ had been placed in any position whatever in the series $r$, $s$ \ldots $l$; and again,
in any position whatever in the series $x$, $y$ \ldots $l$, instead of at the end of each of these.

With a very slight modification this may be made to suit the case of an odd number
instead of one of the two even numbers $2p$, $2q$; (in fact, it is only necessary to place
the mark $(\dagger)$ in $\{AB|\ldots\}$ over the column corresponding to the marked column in
$\{A|\ldots\}$, $\{A|\ldots\}$ being the symbol for which the number of columns is odd), but it is
inapplicable where the two numbers are odd. Consider the second side of (27); this
may be expanded in the form
\[
\Sigma\,\pm_s\,\ldots\pm_x\pm_y\,\cdot\,\overline{AB}_{\rho_1\sigma_1\ldots\,x_1y_1\ldots}\,\cdots\,\overline{AB}_{\rho_k\sigma_k\ldots\,x_ky_k\ldots}
\tag{29},
\]
where $\Sigma$ refers to the different quantities $s,\ldots,x,y,\ldots$ as in (11).

Substituting from (28), this becomes
\[
\Sigma\cdot\Sigma_l\,\ldots\Sigma_l\,(\pm_s\,\ldots\pm_x\pm_y\,A_{\rho_1\sigma_1\ldots l_1}\ldots A_{\rho_k\sigma_k\ldots l_k}\cdot
B_{x_1y_1\ldots l_1}\ldots B_{x_ky_k\ldots l_k})
\tag{30}.
\]
Effecting the summation with respect to $x$, $y$ \ldots this becomes
\[
\Sigma\cdot\Sigma_l\,\ldots\Sigma_l\,\pm_s\,\ldots\,A_{\rho_1\sigma_1\ldots l_1}\ldots A_{\rho_k\sigma_k\ldots l_k}\cdot
\left\{
\begin{array}{c}
\stackrel{\dagger}{B}\,1\,\ldots l_1 \\[2pt]
\hline
k\,k\,\ldots l_k
\end{array}
\right\}
\tag{31},
\]
$\Sigma$ now referring to $s,\ldots$ only. The quantity under the sign $\Sigma$ vanishes if any two
of the quantities $l$ are equal, and in the contrary case, we have
\[
\left\{
\begin{array}{c}
\stackrel{\dagger}{B}\,1\,\ldots l_1 \\[2pt]
\hline
k\,k\,\ldots l_k
\end{array}
\right\}
= \pm_l\,\{B\cdot k\cdot 2q\}
\tag{32},
\]
which reduces the above to
\[
\{B\cdot k\cdot 2q\}\cdot\Sigma\,\pm_s\,\ldots\pm_l\,A_{\rho_1\sigma_1\ldots l_1}\ldots A_{\rho_k\sigma_k\ldots l_k}
\tag{33},
\]
$\Sigma$ referring to the quantities $s,\ldots$, and also to the quantities $l$. And this is evidently
equivalent to
\[
\{A\cdot k\cdot 2p\}\{B\cdot k\cdot 2q\}
\tag{34},
\]
the theorem to be proved. It is obvious that when $p=1$, $q=1$, the equation (27),
coincides with the theorem $(\Theta)$, quoted in the introduction to this paper.

\newpage
\setcounter{section}{12}

% ============================================================
% PAPER 13: ON THE THEORY OF LINEAR TRANSFORMATIONS
% Pages 80--94
% ============================================================

\section{}
\addcontentsline{toc}{section}{13. On the Theory of Linear Transformations}
\begin{center}
\textbf{ON THE THEORY OF LINEAR TRANSFORMATIONS.}
\end{center}
\textit{[From the Cambridge Mathematical Journal, vol.\ IV.\ (1845), pp.\ 193---209.]}

\vspace{6pt}

\textsc{The} following investigations were suggested to me by a very elegant paper on
the same subject, published in the Journal by Mr~Boole. The following remarkable
theorem is there arrived at. If a rational homogeneous function $U$, of the $n^{\text{th}}$ order,
with the $m$ variables $x$, $y$ \ldots, be transformed by linear substitutions into a function
$V$ of the new variables $\xi$, $\eta$\ldots; if, moreover, $\theta U$ expresses the function of the
coefficients of $U$, which, equated to zero, is the result of the elimination of the
variables from the series of equations $\partial_x U=0$, $\partial_y U=0$, \&c., and of course $\theta V$ the
analogous function of the coefficients of $V$: then $\theta V = E^{\alpha}\cdot\theta U$, where $E$ is the
determinant formed by the coefficients of the equations which connect $x$, $y$ \ldots with
$\xi$, $\eta$\ldots, and $\alpha=(m-1)^{n-1}$\footnote{The value of $\alpha$ was left undetermined, but Mr~Boole has since informed me, he was acquainted with
it at the time his paper was written; and has given it in a subsequent paper.}.
In attempting to demonstrate this very beautiful property,
it occurred to me that it might be generalised by considering for the function $U$, not
a homogeneous function of the $n^{\text{th}}$ order between $m$ variables, but one of the same
order, containing $n$ sets of $m$ variables, and the variables of each set entering
linearly. The form which Mr~Boole's theorem thus assumes is $\theta V = E'^{\alpha'}\cdot E''^{\alpha''}\,\ldots E^{(n)\alpha^{(n)}}\cdot\theta U$.
This it was easy to demonstrate would be true, if $\theta U$ satisfied a certain system of
partial differential equations. I imagined at first that these would determine the
function $\theta U$, (supposed, in analogy with Mr~Boole's function, to represent the result of
the elimination of the variables from $\partial_{x_1}U=0$, $\partial_{y_1}U=0$, \ldots $\partial_{x_2}U=0$, \&c.): this I afterwards
found was not the case; and thus I was led to a class of functions, including as
a particular case the function $\theta U$, all of them possessed of the same characteristic
property. The system of partial differential equations was without difficulty replaced
by a more fundamental system of equations, upon which, assumed as definitions, the
theory appears to me naturally to depend; and it is this view of it which I intend
partially to develope in the present paper.

\newpage

I have already employed the notation
\[
\left.
\begin{array}{cccc}
\alpha,&\beta,&\gamma,&\delta,\ldots \\
\alpha',&\beta',&\gamma',&\delta',\ldots \\
\alpha'',&\beta'',&\gamma'',&\delta'',\ldots \\
\vdots&\vdots&\vdots&\vdots
\end{array}
\right\}
\tag{1}
\]
\[
\left.
\begin{array}{cccc}
\alpha,&\beta,&\gamma,&\ldots \\
\alpha',&\beta',&\gamma',&\ldots \\
\alpha'',&\beta'',&\gamma'',&\ldots \\
\vdots&\vdots&\vdots
\end{array}
\right\}
\tag{2},
\]
(where the number of horizontal rows is less than that of vertical ones) to denote
the series of determinants, which can be formed out of the above quantities by selecting any system of vertical
rows; these different determinants not being connected together by the sign $+$, or in
any other manner, but being looked upon as perfectly separate.

The fundamental theorem for the multiplication of determinants gives, applied to
these, the formula
\[
\left.
\begin{array}{cccc}
A,&B,&C,&D,\ldots \\
A',&B',&C',&D',\ldots \\
A'',&B'',&C'',&D'',\ldots \\
\vdots
\end{array}
\right\}
= E
\left.
\begin{array}{cccc}
\alpha,&\beta,&\gamma,&\delta,\ldots \\
\alpha',&\beta',&\gamma',&\delta',\ldots \\
\alpha'',&\beta'',&\gamma'',&\delta'',\ldots \\
\vdots
\end{array}
\right\}
\tag{3},
\]
where
\[
\begin{array}{rcl}
A &=& \lambda\alpha+\lambda'\alpha'+\lambda''\alpha''+\ldots, \\
B &=& \lambda\beta+\lambda'\beta'+\lambda''\beta''+\ldots, \\
A' &=& \mu\alpha+\mu'\alpha'+\mu''\alpha''+\ldots, \\
B' &=& \mu\beta+\mu'\beta'+\mu''\beta''+\ldots, \\
&\&\text{c.}
\end{array}
\tag{4},
\]
\[
E = \left|
\begin{array}{cccc}
\lambda,&\mu,&\ldots \\
\lambda',&\mu',&\ldots \\
\vdots&\vdots
\end{array}
\right|
\tag{5},
\]
and the meaning of the equation is, that the terms on the first side are equal, each
to each, to the terms on the second side.

This preliminary theorem being explained, consider a set of arbitrary coefficients,
represented by the general formula
\[
rst\ldots \tag{6},
\]
in which the number of symbolical letters $r$, $s$, \ldots is $n$, and where each of these is
supposed to assume all integer values, from $1$ to $m$ inclusively.

\newpage

Let
\[
{}_r1a_{t_1\ldots},\ {}_r1'a_{t_1'\ldots},\ \ldots
\tag{7}
\]
represent the whole series, taken in any order, in which the first symbolical letter
is $a$. Similarly,
\[
{}_{r_1}a_{s_1t_1\ldots},\ {}_{r_1'}a_{s_1't_1'\ldots},\ \ldots
\tag{8},
\]
the whole series of those in which the second symbolical letter is $a$, and so on.

Imagine a function $u$ of the coefficients, which is simultaneously of the forms
\[
u = H_p
\left|
\begin{array}{ccc}
{}_r1a_{t_1\ldots},& {}_r1'a_{t_1'\ldots},&\ldots \\
{}_r2a_{t_2\ldots},& {}_r2'a_{t_2'\ldots},&\ldots \\
{}_r3a_{t_3\ldots},& {}_r3'a_{t_3'\ldots},&\ldots \\
\vdots&\vdots
\end{array}
\right|
\]
\[
u = H_p
\left|
\begin{array}{ccc}
{}_{r_1}a_{s_1\ldots},& {}_{r_1'}a_{s_1'\ldots},&\ldots \\
{}_{r_2}a_{s_2\ldots},& {}_{r_2'}a_{s_2'\ldots},&\ldots \\
{}_{r_3}a_{s_3\ldots},& {}_{r_3'}a_{s_3'\ldots},&\ldots \\
\vdots&\vdots
\end{array}
\right|
\quad (A),
\]
\&c.; in which $H_p$ denotes a rational homogeneous function of the order $p$. The
function $H$ is not necessarily supposed to be the same in the above equations, and
in point of fact it will not in general be so. The number of equations is of course $=n$.

The function $u$, whose properties we proceed to investigate, may conveniently be
named a ``Hyperdeterminant.'' Any function satisfying any of the equations $(A)$,
without satisfying all of them, will be an ``Incomplete Hyperdeterminant.'' But,
considering in the first place such as are complete---

Let $rst\ldots$ be a new set of coefficients connected with the former ones by a system
of equations of the form
\[
rst\ldots = \lambda\cdot 1st\ldots + \lambda'\cdot 2st\ldots + \ldots + \lambda^{(m)}\cdot m\,st\ldots
\tag{9},
\]
(where the $r$ in $\lambda^{(r)}\ldots$ is not an exponent, but an affix).

Suppose $\mathfrak{u}$ is the same function of these new coefficients that $u$ was of the former
ones. Then consider the first of the equations $(A)$ and the equation (3), and writing
\[
L = \left|
\begin{array}{cc}
\lambda,&\lambda',\ldots \\
\lambda_1,&\lambda_1',\ldots \\
\vdots&\vdots
\end{array}
\right|
\tag{10},
\]
we have immediately the equation
\[
\mathfrak{u} = L^p\cdot u \tag{11}.
\]

Consider the new set of coefficients
\[
rst\ldots = \mu\cdot r1t\ldots + \mu'\cdot r2t\ldots + \ldots + \mu^{(m)}\cdot r\,m\,t\ldots
\tag{12},
\]
and $\mathfrak{u}$ the analogous function of these; then, from the second of the equations $(A)$
and the equation (3), and writing
\[
M = \left|
\begin{array}{cc}
\mu,&\mu',\ldots \\
\mu_1,&\mu_1',\ldots \\
\vdots&\vdots
\end{array}
\right|
\tag{13},
\]
\[
\mathfrak{u} = M^p\cdot u = L^p M^p\cdot u
\tag{14}.
\]

\newpage

In like manner, considering the new coefficients $rst\ldots$, where
\[
rst\ldots = \nu\cdot rs1\ldots + \nu'\cdot rs2\ldots + \ldots + \nu^{(m)}\cdot rs\,m\,\ldots
\tag{15},
\]
the new function $\mathfrak{u}$, and the quantity $N$ given by
\[
N = \left|
\begin{array}{cc}
\nu,&\nu',\ldots \\
\nu_1,&\nu_1',\ldots \\
\vdots&\vdots
\end{array}
\right|
\tag{16},
\]
we have, as before,
\[
\mathfrak{u} = N^p\cdot u = L^p M^p N^p\cdot u
\tag{17},
\]
or
\[
\mathfrak{u} = L^p M^p N^p\cdot u
\tag{18};
\]
whence, generally, denoting the last result by $u'$,
\[
u' = L^p M^p N^p\,\ldots u
\tag{B, 1}.
\]

Consider now the function
\[
\Sigma\Sigma\Sigma\,\ldots(rst\ldots\,x_r\,y_s\,z_t\,\ldots)
\tag{19},
\]
where the $\Sigma$'s refer successively to $r$, $s$, $t$, \ldots, and denote summations from $1$ to $m$
inclusively. If $u$ be looked upon as a derivative from the above function, we may write
\[
u = \partial_{r}\partial_{s}\partial_{t}\,\ldots\,\Sigma\Sigma\Sigma\,\ldots(rst\ldots\,x_r\,y_s\,z_t\,\ldots)
\tag{20}.
\]
Assume
\[
x_r = M_r^{(1)} y_1 + M_r^{(2)} y_2 + \ldots + M_r^{(m)} y_m
\tag{21}.
\]
It is easy to obtain
\[
\Sigma\Sigma\Sigma\,\ldots(rst\ldots\,x_r\,y_s\,z_t\,\ldots) =
\Sigma\Sigma\Sigma\,\ldots(rst\ldots\,X_r\,Y_s\,Z_t\,\ldots)
\tag{22},
\]
and the formula for $(u)$ becomes
\[
\partial_{x}\partial_{y}\partial_{z}\,\ldots\,\Sigma\Sigma\Sigma\,\ldots(rst\ldots\,x_r\,y_s\,z_t\,\ldots)
= L^p M^p N^p\,\ldots
\partial_{X}\partial_{Y}\partial_{Z}\,\ldots\,\Sigma\Sigma\Sigma\,\ldots(rst\ldots\,X_r\,Y_s\,Z_t\,\ldots)
\tag{B, 2}.
\]

Proceeding to obtain the expression of the coefficients $\mathfrak{r}st\ldots$ in terms of the
coefficients $rst\ldots$, we have
\[
\mathfrak{r}st\ldots = \Sigma\Sigma\,\ldots(\lambda^f\,\mu^g\,\nu^h\,\ldots\,fgh\,\ldots)
\tag{C},
\]
where the $\Sigma$'s refer successively to $f$, $g$, $h$, \ldots, denoting summations from $1$ to $m$
inclusively. Having this equation, it is perhaps as well to retain
\[
u' = L^p M^p N^p\,\ldots u
\tag{B, 1},
\]
instead of $(B, 2)$, that form being principally useful in showing the relation of the
function $u$ to the theory of the transformation of functions.

It may immediately be seen, that in the equations $(B)$, $(C)$ we may, if we please,
omit any number of the marks of variation $(\partial)$, omitting at the same time the
corresponding signs $\Sigma$, and the corresponding factors of the series $L$, $M$, $N$ \ldots

Also, if $u$ be such as only to satisfy some of the equations $(A)$; then, if in the
same formulae we omit the corresponding marks $(\partial)$, summatory signs, and terms of the
series $L$, $M$, $N$ \ldots, the resulting equations are still true.

\newpage

From the formulae $(A)$ we may obtain the partial differential equations
\[
\Sigma\Sigma\,\ldots\,\partial_{r_\alpha st\ldots}\partial_{r_\beta st\ldots}\,\ldots u = 0,\quad\text{or}\quad pu,
\tag{D},
\]
according as $\alpha$ is not equal, or is equal, to $\beta$;
and so on: the summatory signs referring in every case to those of the series $r$, $s$, $t$, \ldots,
which are left variable, and extending from $1$ to $m$ inclusively.

To demonstrate this, consider the general form of $u$, as given by the first of the
equations $(A)$. This is evidently composed of a series of terms, each of the form
\[
c\cdot P\cdot Q\cdot R\,\ldots (p\text{ factors}),
\]
in which
\[
P = \left|
\begin{array}{ccc}
{}_1a_{t_1\ldots},& {}_1'a_{t_1'\ldots},&\ldots\ (m\text{ terms}) \\
{}_2a_{t_2\ldots},& {}_2'a_{t_2'\ldots},&\ldots \\
{}_3a_{t_3\ldots},& {}_3'a_{t_3'\ldots},&\ldots \\
\vdots&\vdots
\end{array}
\right|,
\]
$Q$, $R$, \&c.\ being of the same form; and we have
\[
\partial_{r_\alpha st\ldots} P =
\left|
\begin{array}{ccc}
{}_1a_{t_1\ldots},& {}_1'a_{t_1'\ldots},&\ldots \\
{}_2a_{t_2\ldots},& {}_2'a_{t_2'\ldots},&\ldots \\
\partial_{r_\alpha st\ldots},&\ldots
\end{array}
\right|
+\&\text{c.}+\&\text{c.},
\]
and
\[
\Sigma\Sigma\,\ldots\,\partial_{r_\alpha st\ldots}\partial_{r_\beta st\ldots} P =
\left|
\begin{array}{ccc}
{}_1a_{t_1'\ldots},& {}_1'a_{t_1''\ldots},&\ldots\ (m\text{ terms}) \\
{}_2a_{t_2'\ldots},& {}_2'a_{t_2''\ldots},&\ldots \\
\partial_{r_\alpha st\ldots},&\ldots
\end{array}
\right| = 0;
\]
so that all the terms on the second side of the equation vanish. If, however, $\beta=\alpha$,
whence, on the second side, we have
\[
cQR\,\ldots P + cPR\,\ldots Q + \&\text{c.}
= p\cdot cPQR\,\ldots + \&\text{c.}+\&\text{c.} = pu,
\]
or the theorem in question is proved.

In the case of an incomplete hyperdeterminant, the corresponding systems of
equations are of course to be omitted. In every case it is from these equations that
the form of the function $u$ is to be investigated; they entirely replace the system $(A)$.

A very important case of the general theory is, when we suppose the coefficients
$rst\ldots$ to have the property $r's't'\ldots=r''s''t''\ldots$, whenever $r's't'\ldots$ and $r''s''t''\ldots$ denote
the same combination of letters; and also that the coefficients $\lambda$ are equal to the
coefficients $\mu$, $\nu$ \ldots, each to each. In this case the coefficients $\mathfrak{r}st\ldots$ have likewise
the same property, viz.\ that $\mathfrak{r}''s''t''\ldots = \mathfrak{r}'s't'\ldots$, whenever $r's't'\ldots$ and $r''s''t''\ldots$ denote
the same combination of letters.

\newpage

The equations $(B, 1)$, $(B, 2)$, become in this case
\[
u' = L^{np}\cdot u
\tag{B, 3},
\]
instead of $(B, 1)$, that form being principally useful in showing the relation of the
function $u$ to the theory of the transformation of functions.

It may immediately be seen, that in the equations $(B)$, $(C)$ we may, if we please,
omit any number of the marks of variation $(\partial)$, omitting at the same time the
corresponding signs $\Sigma$, and the corresponding factors of the series $L$, $M$, $N$ \ldots

Also, if $u$ be such as only to satisfy some of the equations $(A)$; then, if in the
same formulae we omit the corresponding marks $(\partial)$, summatory signs, and terms of the
series $L$, $M$, $N$ \ldots, the resulting equations are still true.

so that all the terms on the second side of the equation vanish. If, however, $\beta=\alpha$,
\[
\Sigma\Sigma\,\ldots\,\left(\mathrm{ast}\ldots\,\dfrac{d}{d\beta st\ldots}\right)P = P;
\]
whence, on the second side of the equation vanish. If, however,
\[
cQR\ldots P + cPR\ldots Q + \&\text{c.} = p\cdot cPQR\ldots = pU,
\]
or the theorem in question is proved.

In the notation there employed, we have
\[
u =
\left\{
\begin{array}{c}
\stackrel{\dagger}{1}\,\stackrel{\dagger}{1}\,\ldots(n) \\[2pt]
\hline
2\,2 \\
\vdots \\
mm
\end{array}
\right\},
\]
a complete hyperdeterminant when $n$ is even; and when $n$ is odd the functions
\[
\left\{
\begin{array}{c}
\stackrel{\dagger}{1}\,1\,\ldots(n) \\[2pt]
\hline
2\,2 \\
\vdots \\
mm
\end{array}
\right\},
\quad
\left\{
\begin{array}{c}
1\,\stackrel{\dagger}{1}\,\ldots(n) \\[2pt]
\hline
2\,2 \\
\vdots \\
mm
\end{array}
\right\}
\]
are each of them incomplete hyperdeterminants.

$(A)$. In the case of $n=2$, the complete hyperdeterminant is simply the ordinary
determinant
\[
\left|
\begin{array}{cccc}
1\,1,&1\,2,&\ldots&1\,m \\
2\,1,&2\,2,&\ldots&2\,m \\
\vdots&\vdots&&\vdots \\
m\,1,&m\,2,&\ldots&mm
\end{array}
\right|.
\]
Stating the general conclusion as applied to this case, which is a very well known one,
``If the function
\[
U = \Sigma\Sigma\,(rs\cdot x_ry_s),
\]
be transformed into a similar function
\[
\Sigma\Sigma\,(\mathfrak{rs}\cdot\bar{x}_r\bar{y}_s),
\]
by means of the substitutions
\[
x_r = \lambda_1^r\,\bar{x}_1 + \lambda_2^r\,\bar{x}_2 + \ldots + \lambda_m^r\,\bar{x}_m,
\]
\[
y_s = \mu_1^s\,\bar{y}_1 + \mu_2^s\,\bar{y}_2 + \ldots + \mu_m^s\,\bar{y}_s,
\]
then
\[
\left|
\begin{array}{ccc}
1\,1,&1\,2,&\ldots \\
2\,1,&2\,2,&\ldots \\
\vdots
\end{array}
\right|
=
\left|
\begin{array}{ccc}
\lambda_1^1,&\lambda_2^1,&\ldots \\
\lambda_1^2,&\lambda_2^2,&\ldots \\
\vdots
\end{array}
\right|
\left|
\begin{array}{ccc}
\mu_1^1,&\mu_2^1,&\ldots \\
\mu_1^2,&\mu_2^2,&\ldots \\
\vdots
\end{array}
\right|
\left|
\begin{array}{ccc}
\mathfrak{1}\mathfrak{1},&\mathfrak{1}\mathfrak{2},&\ldots \\
\mathfrak{2}\mathfrak{1},&\mathfrak{2}\mathfrak{2},&\ldots \\
\vdots
\end{array}
\right|.
\]
Also, by what has preceded,
\[
\mathfrak{rs} = \Sigma\Sigma\,(\lambda_f^r\cdot\mu_g^s\cdot fg);
\]
so that the theorem is easily seen to amount to the following one:---``If the terms of
a determinant of the $m^{\text{th}}$ order be of the form $\Sigma\Sigma_g(rs,x_r,y_g)$, $x$, $y$ extending
before, from 1 to $m$ inclusively, the determinant itself is the product of three deter-
minants; the first formed with the coefficients $rs$, the second with the quantities $x$,
and the third with the quantities $y$.''

\newpage

In a following number of the \textit{Journal} I shall prove, and apply to the theories of
Maxima and Minima and of Spherical Coordinates, (I may just mention having obtained,
in an elegant form, the formulae for transforming from one oblique set of coordinates
to another oblique one) the more general theorem.

``If $k$ be the order of the determinant formed as above, the determinant itself
is a quadratic function, its coefficients being determinants formed with the coefficients
$rs$, its variables being determinants formed respectively with the variables $x$ and the
variables $y$; and the number of variables in each set being the number of combi-
nations of $k$ things out of $m$, $(=1$ if $k=m$; if $k>m$ the determinant vanishes)\footnote{See concluding paragraph of this paper.}.''

I shall give in the same paper the demonstration of a very beautiful theorem,
rather relating, however, to determinants than to quadratic functions, proved by Hesse
in a Memoir in \textit{Crelle's Journal}, vol.\ XX.\ ``De curvis et superficiebus secundi ordinis;''
and from which he has deduced the most interesting geometrical results. Another
Memoir, by the same author, \textit{Crelle}, vol.\ XXVIII.\ ``Ueber die Elimination der Variabeln
aus drei algebraischen Gleichungen vom zweiter Grade mit zwei Variabeln,'' though
relating in point of fact rather to functions of the third order, contains some most
important results. A few theorems on quadratic functions, belonging, however, to a
different part of the subject, will be found in my paper already quoted in the \textit{Cambridge
Philosophical Transactions} [12]; and likewise in a paper in the \textit{Journal}, Chapters in
the Algebraical Geometry of $n$ dimensions [11].

I shall, just before concluding this case, write down the particular formula corre-
sponding to three variables, and for the symmetrical case. It is, as is well known, the
theorem.

``If
\[
U = Ax^2 + By^2 + Cz^2 + 2Fyz + 2Gxz + 2Hxy
\]
be transformed into
\[
\mathfrak{A}\xi^2 + \mathfrak{B}\eta^2 + \mathfrak{C}\zeta^2 + 2\mathfrak{F}\eta\zeta + 2\mathfrak{G}\xi\zeta + 2\mathfrak{H}\xi\eta
\]
by means of
\[
x = \alpha\xi + \beta\eta + \gamma\zeta,
\]
\[
y = \alpha'\xi + \beta'\eta + \gamma'\zeta,
\]
\[
z = \alpha''\xi + \beta''\eta + \gamma''\zeta,
\]
then
\[
(\mathfrak{A}\mathfrak{B}\mathfrak{C} - \mathfrak{A}\mathfrak{F}^2 - \mathfrak{B}\mathfrak{G}^2 - \mathfrak{C}\mathfrak{H}^2 + 2\mathfrak{F}\mathfrak{G}\mathfrak{H}) =
\]
\[
(\alpha\beta'\gamma'' - \alpha\beta''\gamma' + \alpha'\beta''\gamma - \alpha'\beta\gamma'' + \alpha''\beta\gamma' - \alpha''\beta'\gamma)^2
(ABC - AF^2 - BG^2 - CH^2 + 2FGH).''
\]

$(B)$ Let $n=3$, and for greater simplicity $m=2$; write
\[
\begin{array}{cc}
a = 111,& e = 112, \\
b = 211,& f = 212, \\
c = 121,& g = 122, \\
d = 221,& h = 222,
\end{array}
\]
so that
\[
U = a\,x_1y_1z_1 + b\,x_2y_1z_1 + c\,x_1y_2z_1 + d\,x_2y_2z_1 + e\,x_1y_1z_2 + f\,x_2y_1z_2 + g\,x_1y_2z_2 + h\,x_2y_2z_2.
\]

\newpage

There is no complete hyperdeterminant (i.e.\ for $p=1$), and the incomplete ones are
\begin{align*}
ah-bg-cf+de &= u, \\
ah-de-bg+cf &= u_1, \\
ah-cf-de+bg &= u_2,
\end{align*}
suppose. Thus, suppose the transforming equations are
\[
x_1 = \lambda_1^1\,\bar{x}_1 + \lambda_2^1\,\bar{x}_2,\quad
x_2 = \lambda_1^2\,\bar{x}_1 + \lambda_2^2\,\bar{x}_2,
\]
\[
y_1 = \mu_1^1\,\bar{y}_1 + \mu_2^1\,\bar{y}_2,\quad
y_2 = \mu_1^2\,\bar{y}_1 + \mu_2^2\,\bar{y}_2,
\]
\[
z_1 = \nu_1^1\,\bar{z}_1 + \nu_2^1\,\bar{z}_2,\quad
z_2 = \nu_1^2\,\bar{z}_1 + \nu_2^2\,\bar{z}_2;
\]
then
\[
\bar{u}_1 = (\lambda_1^1\lambda_2^2 - \lambda_2^1\lambda_1^2)(\mu_1^1\mu_2^2 - \mu_2^1\mu_1^2)(\nu_1^1\nu_2^2 - \nu_2^1\nu_1^2)\,u_1,\quad
\text{where $y$, $z$ are changed.}
\]
\[
\bar{u}_2 = (\lambda_1^1\lambda_2^2 - \lambda_2^1\lambda_1^2)(\mu_1^1\mu_2^2 - \mu_2^1\mu_1^2)(\nu_1^1\nu_2^2 - \nu_2^1\nu_1^2)\,u_2,\quad
\text{$z$, $x$.}
\]
\[
\bar{u}\phantom{_1} = (\lambda_1^1\lambda_2^2 - \lambda_2^1\lambda_1^2)(\mu_1^1\mu_2^2 - \mu_2^1\mu_1^2)(\nu_1^1\nu_2^2 - \nu_2^1\nu_1^2)\,u,\quad
\text{$x$, $y$.}
\]

We might also have assumed
\begin{align*}
u_1 &= ad - bc,\quad eh - gf, \\
u_2 &= af - be,\quad ch - dg, \\
u_3 &= ag - ce,\quad bh - df,
\end{align*}
but these are ordinary determinants.

$(C)$ $n=4$, $m=2$. Assume
\[
\begin{array}{cccc}
a = 1111,& i = 1112,& & \\
b = 2111,& j = 2112,& & \\
c = 1211,& k = 1212,& & \\
d = 2211,& l = 2212,& & \\
e = 1121,& m = 1122,& & \\
f = 2121,& n = 2122,& & \\
g = 2211,& o = 2212,& & \\
h = 2221,& p = 2222.& &
\end{array}
\]
\[
\begin{aligned}
U &= a\,x_1y_1z_1w_1 + b\,x_2y_1z_1w_1 + c\,x_1y_2z_1w_1 + d\,x_2y_2z_1w_1 \\
&\quad + e\,x_1y_1z_2w_1 + f\,x_2y_1z_2w_1 + g\,x_1y_2z_2w_1 + h\,x_2y_2z_2w_1 \\
&\quad + i\,x_1y_1z_1w_2 + j\,x_2y_1z_1w_2 + k\,x_1y_2z_1w_2 + l\,x_2y_2z_1w_2 \\
&\quad + m\,x_1y_1z_2w_2 + n\,x_2y_1z_2w_2 + o\,x_1y_2z_2w_2 + p\,x_2y_2z_2w_2,
\end{aligned}
\]
we have
\[
u = ap - bo - cn + dm - el + fk + gj - hi.
\]

\newpage

so that, with the same sets of transforming equations as above, and the additional one,
\[
w_1 = \rho_1^1\,\bar{w}_1 + \rho_2^1\,\bar{w}_2,
\]
\[
w_2 = \rho_1^2\,\bar{w}_1 + \rho_2^2\,\bar{w}_2,
\]
we have
\[
\bar{u} = (\lambda_1^1\lambda_2^2 - \lambda_2^1\lambda_1^2)(\mu_1^1\mu_2^2 - \mu_2^1\mu_1^2)(\nu_1^1\nu_2^2 - \nu_2^1\nu_1^2)(\rho_1^1\rho_2^2 - \rho_2^1\rho_1^2)\,u;
\]
this is important when viewed in reference to a result which will presently be obtained.

If we take the symmetrical case, we have
\[
U = ax^4 + 4\beta x^3y + 6\gamma x^2y^2 + 4\delta xy^3 + \epsilon y^4;
\]
which is transformed into
\[
U = \mathfrak{a}\bar{x}^4 + 4\mathfrak{b}\bar{x}^3\bar{y} + 6\mathfrak{c}\bar{x}^2\bar{y}^2 + 4\mathfrak{d}\bar{x}\bar{y}^3 + \mathfrak{e}\bar{y}^4;
\]
by means of
\[
x = \lambda\bar{x} + \mu\bar{y},
\]
\[
y = \lambda'\bar{x} + \mu'\bar{y};
\]
then, if
\[
u = \alpha\epsilon - 4\beta\delta + 3\gamma^2,
\]
\[
u' = \mathfrak{a}\mathfrak{e} - 4\mathfrak{b}\mathfrak{d} + 3\mathfrak{c}^2,
\]
we have
\[
u' = (\lambda\mu' - \lambda'\mu)^4\cdot u.
\]

\noindent\textbf{II.} Where $p=2$, $m=2$, $n$ is odd.

The expression
\[
u =
\left\{
\begin{array}{c}
\stackrel{\dagger}{1}\,1\,1\,\ldots(n) \\[2pt]
\hline
2\,2\,2\,\ldots
\end{array}
\right\}^2,
\quad
\left\{
\begin{array}{c}
\stackrel{\dagger}{1}\,1\,1\,1\,\ldots(n) \\[2pt]
\hline
2\,2\,2\,2\,\ldots
\end{array}
\right\}^2,
\quad
\left\{
\begin{array}{c}
1\,\stackrel{\dagger}{1}\,1\,\ldots(n) \\[2pt]
\hline
2\,2\,2\,\ldots
\end{array}
\right\}^2,
\quad
\left\{
\begin{array}{c}
1\,1\,\stackrel{\dagger}{1}\,\ldots(n) \\[2pt]
\hline
2\,2\,2\,\ldots
\end{array}
\right\}^2
\]
is a complete hyperdeterminant; and that over whichever the mark $(\dagger)$ of nonper-
mutation is placed. The different expressions so obtained are not, however, all of them
independent functions: thus, in the following example, where $n=3$, the three functions
are absolutely identical.

$(A)$. $n=3$, notation as in I.\ $(B)$.
\[
u = a^2h^2 - b^2g^2 + c^2f^2 + d^2e^2 - 2abhg - 2acfh - 2bdeh - 2cdef + 4adfg + 4bceh,
\]
and then
\[
\bar{u} = (\lambda_1^1\lambda_2^2 - \lambda_2^1\lambda_1^2)^2(\mu_1^1\mu_2^2 - \mu_2^1\mu_1^2)^2(\nu_1^1\nu_2^2 - \nu_2^1\nu_1^2)^2\,u.
\]

This is in many respects an interesting example. We see that the function $u$
may be expressed in the three following forms:
\begin{align*}
u &= (ah - bg - cf + de)^2 + 4(ad - bc)(fg - eh), \tag{1}\\
u &= (ah - bg - de + cf)^2 + 4(af - be)(dg - ch), \tag{2}\\
u &= (ah - cf - de + bg)^2 + 4(ag - ce)(df - bh). \tag{3}
\end{align*}

\newpage

(Pages 91--93 contain a large symbolic expansion of the function for the case
$n=4$, $m=2$, given as a list of monomials in the coefficients $a,b,c,d,e,f,g,h,i,j,k,l,m,n,o,p$
and constants $\mathfrak{A,B,C,D,E,F,G,H}$. The detailed expansion is reproduced in
Cayley's original; here we give the leading expressions in symbolic form.)

We have, as usual,
\[
\bar{u} = (\lambda_1^1\lambda_2^2 - \lambda_2^1\lambda_1^2)^2(\mu_1^1\mu_2^2 - \mu_2^1\mu_1^2)^2(\nu_1^1\nu_2^2 - \nu_2^1\nu_1^2)^2(\rho_1^1\rho_2^2 - \rho_2^1\rho_1^2)^2\,u.
\]

Particular forms of $u$ are
\[
A=1,\quad B=0:
\]
\[
u = (\mathfrak{A} + 3\mathfrak{B} + 3\mathfrak{C} + 6\mathfrak{D}) = (ap - bo - cn + dm - el + fk + gj - hi) = \nu \quad\text{suppose.}
\]
\[
A=1,\quad B=9:
\]
\[
u = (\mathfrak{A} + 3\mathfrak{B} - 6\mathfrak{C} - 5\mathfrak{D} + 33\mathfrak{E} + 9\mathfrak{F} - 18\mathfrak{G} - 27\mathfrak{H}) = \theta U \quad\text{suppose,}
\]
where $\theta U=0$ is the result of the elimination of the variables from the equations
$\partial_x U=0$, $\partial_y U=0$, $\partial_z U=0$, $\partial_w U=0$, $\partial_{x'}U=0$, $\partial_{y'}U=0$, $\partial_{z'}U=0$, $\partial_{w'}U=0$. In fact, by
an investigation similar to Mr~Boole's, applied to a function such as $U$, it shews
that $\theta U$ has the characteristic property of the function $u$; also in the present case
$u$ is the most general function of this kind, so that $\theta U$ is obtained from $U$ by
properly determining the constant. This has been effected by comparing the value of
$u$, in the symmetrical case, with the value of $\theta U$, in the same case, the expanded
expression of which is given by Mr Boole in the \textit{Journal}, vol.\ IV.\ p.\ 169. Assuming
$A=1$, the result was $B=9$. [Incorrect: the result of the elimination is not $\theta U=0$,
but an equation of a higher degree.]

The general form of $u$ now becomes
\[
u = ax^2 + \beta\bar{U},
\]
in which $a$, $\beta$, are indeterminate.

We have
\[
\bar{u} = a\bar{x}^2 + \beta\bar{\bar{U}} = M\,(a x^2 + \beta\bar{U}),
\]
where
\[
M = (\lambda_1^1\lambda_2^2 - \lambda_2^1\lambda_1^2)^2(\mu_1^1\mu_2^2 - \mu_2^1\mu_1^2)^2(\nu_1^1\nu_2^2 - \nu_2^1\nu_1^2)^2(\rho_1^1\rho_2^2 - \rho_2^1\rho_1^2)^2,
\]
and thence
\[
\bar{\nu}^2 = M\nu^2,
\]
which coincides with a previous formula, and
\[
\theta\bar{U} = M\theta U;
\]
whence, eliminating $M$,
\[
\frac{\theta\bar{U}}{\bar{\nu}^2} = \frac{\theta U}{\nu^2};
\]
an equation which is remarkable as containing only the constants of $U$ and $\bar{U}$; it is
an equation of condition which must exist among the constants of $\bar{U}$ in order that
this function may be derivable by linear substitutions from $U$.

\newpage

In the symmetrical case, or where
\[
U = ax^4 + 4\beta x^3y + 6\gamma x^2y^2 + 4\delta xy^3 + \epsilon y^4,
\]
it has been already seen that $\nu$ is given by
\[
\nu = \alpha\epsilon - 4\beta\delta + 3\gamma^2.
\]
Proceeding to form $\theta U$, we have
\begin{align*}
\mathfrak{A} &= \alpha\epsilon^3 - 4\beta\delta^3, \\
\mathfrak{B} &= 4(\alpha\delta^2\epsilon - \alpha\epsilon^2\delta + 3\beta\gamma\epsilon - 3\beta\delta^2\gamma), \\
\mathfrak{C} &= 3(\alpha^2\epsilon^2 + 2\alpha\gamma\epsilon^2 - 4\alpha\delta^2\epsilon), \\
\mathfrak{D} &= 6(\alpha\gamma\epsilon^2 - 2\alpha\delta^2\epsilon + 3\delta^2\gamma\epsilon - 2\delta^2\gamma^2), \\
\mathfrak{E} &= 5\alpha\gamma^3 - 4\beta^3\gamma + \&\text{c.}, \\
\mathfrak{F} &= 6(\alpha^2\gamma\epsilon - 2\alpha\beta\delta\epsilon - \alpha\gamma^2\delta + 4\beta^2\gamma\delta - 4\beta\gamma^3 + \&\text{c.}), \\
\mathfrak{G} &= 12(\alpha\beta\delta\gamma - \alpha\beta\delta^2 - \alpha\gamma^2\delta + 2\gamma^3\delta + \beta^2\gamma\delta - 2\beta^3\delta), \\
\mathfrak{H} &= 6(\alpha^2\delta^3 - 3\alpha^2\delta^2\epsilon + 4\alpha\beta\delta\epsilon - 6\beta^2\gamma\delta + \&\text{c.}),
\end{align*}
and these values give
\[
\theta U = \alpha^2\epsilon^3 - 6\alpha\beta^2\epsilon^2 - 12\alpha\delta^2\beta\epsilon + 18\alpha\gamma^2\beta^2 + 27\beta^2\delta^4 + 56\delta^2\beta^4\epsilon
- 24\beta^4\gamma^3 + \&\text{c.}
\]
so that this function, divided by $(\alpha\epsilon - 4\beta\delta + 3\gamma^2)^2$, is invariable for all functions of the
fourth order which can be deduced one from the other by linear substitutions. The
function $u = 4\beta^2 - 3\gamma^2$ occurs in other investigations; I have met with it in a problem
relating to a homogeneous function of two variables, of any order whatever, $\alpha$, $\beta$, $\gamma$, $\delta$, $\epsilon$
signifying the fourth differential coefficients of the function. But this is only remotely
connected with the present subject.

Since writing the above, Mr~Boole has pointed out to me that in the transform-
ation of a function of the fourth order of the form $\alpha x^4 + 4\beta x^3y + 6\gamma x^2y^2 + 4\delta xy^3 + \epsilon y^4$,---
besides his function $\theta v$, and my quadratic function $\alpha\epsilon - 4\beta\delta + 3\gamma^2$,---there exists a function

\newpage

of the third order $\alpha\epsilon - b^2 - \alpha c^2 - c^3 + 2bcd$, possessing precisely the same characteristic
property, and that, moreover, the function $\theta v$ may be reduced to the form
\[
(\alpha\epsilon - 4\beta\delta + 3\gamma^2)^3 - 27(\alpha\epsilon - \alpha\delta^2 - c^3 - c^2 + 2bcd)^2;
\]
the latter part of which was verified by trial; the former he has demonstrated in a
manner which, though very elegant, does not appear to be the most direct which Mr
Boole, in his earliest paper on the subject, \textit{Mathematical Journal}, vol.\ II.\ p.\ 70.
The equations $d_xv=0$, $d_yv=0$, $d_zv=0$, imply the corresponding equations for the
transformed function: from these equations we might obtain two relations between the
coefficients, which, in the case of a function of the fourth order, are of the orders 3
and 4 respectively: these imply the corresponding relations between the coefficients of
the transformed function. Let $A=0$, $B=0$, $A'=0$, $B'=0$, must have $A'=AA+MB$, $A$, $A'$ being
of the order 3 in the coefficients of $v$, but $B$, $B'$ only of
the third order in the coefficients of $v$, it is evident that the term $MB$ must disappear,
or that the equation is of the form $A'=\lambda A$. The function $A$ is obviously the function
which, equated to zero, would imply the elimination of $\alpha$, $\beta$, $\gamma$, $\delta$, considered
as independent quantities from the equations $\alpha x^4 + 4\beta x^3y + \&\text{c.} = 0$, considered as
linear in $\lambda$, $\mu$, but $\lambda$, $\mu'$ being arbitrary, this implies the corresponding equations of
the third order, would lead to the result $A=AA$. The function $A$ is obviously the
fourth order, for functions of the sixth order, I have actually expanded: but it does not
appear to contain the complete theory. Again, in the particular case of homogeneous
functions of two variables, the transforming functions may be expressed as symmetrical
functions of the roots of the equation $u=0$, which gives rise to an entirely distinct
theory. This, however, I have not as yet developed sufficiently for publication. There
does not appear to be anything very directly analogous to the subject of this note, in
my general theory: if this be so, it proves the absolute necessity of a distinct investi-
gation for the present case, that which I have denominated the symmetrical one.

\newpage

of the third order $\alpha\epsilon - b\eta - c\theta^2 + 2bcde$, possessing precisely the same characteristic
property, and that, moreover, the function $\theta v$ may be reduced to the form
\[
(\alpha\epsilon - 4bd + 3c^2)^3 - 27(\alpha ce - \alpha d^2 - eb^2 - c^3 + 2bcd)^2;
\]
the latter part of which was verified by trial; the former he has demonstrated in a
manner which, though very elegant, does not appear to be the most direct which can
be obtained. The corresponding functions, which in the transformation of $\partial v$ may be
considered: from these equations we might obtain two relations between the
coefficients, which, in the case of a function of the fourth order, are of the orders 3
and 4 respectively: these imply the corresponding relations between the coefficients of
the transformed function. Let $A=0$, $B=0$, $A'=0$, $B'=0$, imply the corresponding equations
that $A$, $\lambda$, $\mu$, but $\lambda$, $\mu'$ being arbitrary, this implies the corresponding equations of
the third order, would lead to the result $A=AA$. The function $A$ is obviously the
function which, equated to zero, would express the elimination of $\alpha$, $\beta$, $\gamma$, $\delta$,
considered as independent quantities from the equations $\alpha x^4 + 4\beta x^3y + \&\text{c.} = 0$, considered as
linear in $\lambda$, $\mu$, but $\lambda$, $\mu'$ being arbitrary, this implies the corresponding equations of
the third order, would lead to the result $A=AA$. The function $A$ is obviously the
fourth order, for functions of the sixth order, I have actually expanded: but it does not
appear to contain the complete theory. Again, in the particular case of homogeneous
functions of two variables, the transforming functions may be expressed as symmetrical
functions of the roots of the equation $u=0$, which gives rise to an entirely distinct
theory. This, however, I have not as yet developed sufficiently for publication. There
does not appear to be anything very directly analogous to the subject of this note, in
my general theory: if this be so, it proves the absolute necessity of a distinct investi-
gation for the present case, that which I have denominated the symmetrical one.

\newpage
\setcounter{section}{13}

% ============================================================
% PAPER 14: ON LINEAR TRANSFORMATIONS
% Pages 95--112 (this chunk: pp. 95-100)
% ============================================================

\section{}
\addcontentsline{toc}{section}{14. On Linear Transformations}
\begin{center}
\textbf{ON LINEAR TRANSFORMATIONS.}
\end{center}
\textit{[From the Cambridge and Dublin Mathematical Journal, vol.\ I.\ (1846), pp.\ 104---122.]}

\vspace{6pt}

\textsc{In} continuing my researches on the present subject, I have been led to a new
manner of considering the question, which, at the same time that it is much more
general, has the advantage of applying directly to the only case which one can
possibly hope to develope with any degree of completeness, that of functions of two
variables. In fact the question may be proposed, ``To find all the derivatives of any
number of functions, which have the property of preserving their form unaltered after
any linear transformations of the variables.'' By Derivative I understand a function
deduced in any manner whatever from the given functions, and I give the name of
Hyperdeterminant Derivative, or simply of Hyperdeterminant, to those derivatives which
have the property just enunciated. These derivatives may easily be expressed explicitly,
by means of the known method of the separation of symbols. We thus obtain the
most general expression of a hyperdeterminant. But there remains a question to be
resolved, which appears to present very great difficulties, that of determining the
independent derivatives, and the relation between these and the remaining ones. I
have only succeeded in treating a very particular case of this question, which shows
however in what way the general problem is to be attacked.

Imagine $p$ series each of $m$ variables
\[
x_1,\ y_1,\ldots\&\text{c.}\quad x_2,\ y_2,\ldots\&\text{c.}\quad x_p,\ y_p,\ldots\&\text{c.},
\]
where $p$ is at least as great as $m$.

Similarly $p'$ series each of $m'$ variables
\[
x_1',\ y_1',\ldots\&\text{c.}\quad x_2',\ y_2',\ldots\&\text{c.}\quad \ldots x_{p'}',\ y_{p'}',\ldots\&\text{c.},
\]

\newpage

$p'$ at least as great as $m'$, and so on. Let the analogous variables $\bar{x}$, $\bar{y}$, \ldots be con-
nected with these by the equations
\[
x = \lambda\bar{x} + \mu\bar{y} + \ldots,
\]
\[
y = \lambda'\bar{x} + \mu'\bar{y} + \ldots,
\]
\[
\vdots
\]
\[
x' = \lambda^1\bar{x}' + \mu^1\bar{y}' + \ldots,
\]
\[
y' = \lambda'^1\bar{x}' + \mu'^1\bar{y}' + \ldots,
\]
\[
\vdots
\]
where $x$, $y$, \ldots\ or $x_1$, $y_1$, \ldots\ or $x_p$, $y_p$, \ldots\ stand for $x_1$, $y_1$, \ldots; $x_p'$, $y_p'$, \ldots\ stand for $x_1'$, $y_1'$, \ldots,
or $x_{p'}'$, $y_{p'}'$, \ldots\ or $x_{p'}'$, $y_{p'}'$, \&c.\ldots\ The coefficients $\lambda$, $\mu$, \ldots, $\lambda'$, $\mu'$, \ldots, \&c.; $\lambda^1$, $\mu^1$, \ldots, $\lambda'^1$, $\mu'^1$, \ldots
remain the same in all these systems. Suppose next,
\[
\xi = \delta_x,\quad \eta = \delta_y,
\]
i.e.
\[
\xi_1 = \delta_{x_1},\quad \eta_1 = \delta_{y_1},\ldots, \xi_p = \delta_{x_p},\ldots
\]
(where $\delta_x$, $\delta_y$ \ldots are the symbols of differentiation relative to $x$, $y$, \&c.). Then evidently
\[
\xi = \mu'\bar{\xi} + \lambda'\bar{\eta} + \ldots,
\]
\[
\eta = \mu\bar{\xi} + \lambda\bar{\eta} + \ldots,
\]
\[
\vdots
\]
with similar equations for $\xi'$, $\eta'$, \ldots Suppose
\[
\|\Omega\| =
\left|
\begin{array}{ccc}
\xi_1,&\xi_2,&\ldots\xi_p \\
\eta_1,&\eta_2,&\ldots\eta_p \\
\vdots
\end{array}
\right|,
\quad
\|\Omega'\| =
\left|
\begin{array}{ccc}
\xi_1',&\xi_2',&\ldots\xi_{p'}' \\
\eta_1',&\eta_2',&\ldots\eta_{p'}' \\
\vdots
\end{array}
\right|,\ldots
\]
that is to say $\|\Omega\|$ is the series of determinants formed by choosing any $m$ vertical
columns to compose a determinant, and similarly $\|\Omega'\|$, \&c. Suppose, besides,
\[
E =
\left|
\begin{array}{ccc}
\lambda,&\mu,&\ldots \\
\lambda',&\mu',&\ldots \\
\vdots
\end{array}
\right|,
\quad
E' =
\left|
\begin{array}{ccc}
\lambda^1,&\mu^1,&\ldots \\
\lambda'^1,&\mu'^1,&\ldots \\
\vdots
\end{array}
\right|,\ldots
\]
Then, by the known properties of determinants,
\[
\|\bar{\Omega}\| = E\|\Omega\|,\quad \|\bar{\Omega}'\| = E'\|\Omega'\|,\ \&\text{c.,}
\]
i.e.\ the terms on the one side are respectively equal to the terms on the other. Hence if
\[
\square = F(\|\Omega\|,\ \|\Omega'\|^{f'},\ \ldots),
\]
i.e.\ $\square$ a rational and integral function, homogeneous of the order $f$ in the quantities
of the series $\|\Omega\|$, homogeneous of the order $f'$ in the quantities of the series $\|\Omega'\|$,
\&c., we have immediately
\[
\bar{\square} = E^f E'^{f'}\ldots\square;
\]

\newpage

or if $U$ be any function whatever of the variables $x$, $y$\ldots\ which is transformed by
the linear substitutions above into $\bar{U}$, then
\[
\bar{\square}\,\bar{U} = E^f E'^{f'}\ldots\square U,
\]
or the function
\[
\square U
\]
is by the above definition a hyperdeterminant derivative. The symbol $\square$ may be called
``symbol of hyperdeterminant derivation,'' or simply ``hyperdeterminant symbol.''

Let $A$, $B$,\ldots represent the different quantities of the series $\|\Omega\|$,---$A'$, $B'$,\ldots those of
the series $\|\Omega'\|$, \&c.,\ldots, then $\square$ may be reduced to any single term, and we may write
\[
\square = A^\alpha B^\beta\ldots A'^{\alpha'} B'^{\beta'}\ldots
\]
Also $U$ may be supposed of the form
\[
U = \Theta\Phi\ldots
\]
where $\Theta$, $\Phi$ are functions of the variables of one of the sets $x$, $y$, \ldots, of one of the
sets $x'$, $y'$, \&c., thus $\Theta$ is of the form
\[
F(x,\ y_1,\ldots\,x_2,\ y_2,\ldots\,x_p,\ y_p,\ldots),
\]
and so on. The functions $\Theta$, $\Phi$, \ldots may be the same or different. It may be supposed
after the differentiations that several of the sets $x$, $y$, \ldots, or of the sets $x'$, $y'$, \ldots,
become identical; in such cases it will always be assumed that the functions $\Theta$, \ldots
into the differentiations enter, are similar; and the differential coefficients absolutely
identical, in such cases the variables they contain are made so. Thus the general expression
of a hyperdeterminant is placed. The different expressions so obtained are not, however, all of them
independent functions: thus, in the following example, where $n=3$, the three functions
are absolutely identical.

What precedes, in the general theory; it might, perhaps have been more clearly
by confining it to a particular case: and by doing this from the beginning; it will
be seen that it presents no real difficulties. Passing at present to some developments,
to do this, I neglect entirely the sets $x'$, $y'$\ldots\ and I assume that the number $m$ of
variables in each of the sets $x$, $y$, \ldots\ and I assume that the number $m$ of
variables in each of the sets $x$, $y$, \ldots is equal to two. I consider therefore functions
$V_1$, $V_2$, $V_3$ of the variables $x_1$, $y_1$, or $x_2$, $y_2$, \ldots\ or $x_p$, $y_p$. Writing also
\[
\xi_r\eta_s - \xi_s\eta_r = \overline{12},\ \&\text{c.}
\]
the symbols $A$, $B$, \ldots reduce themselves to $\overline{12}$, $\overline{13}$, \ldots Hence for functions of two
variables, there results the following still tolerably general form
\[
\square = \overline{12}^a\overline{13}^b\overline{14}^c\ldots\overline{23}^{a'}\overline{24}^{b'}\overline{34}^{a''}\ldots\,V_1V_2V_3V_4\ldots
\]

\newpage

The functions $V_1$, $V_2$,\ldots may be the same or different; but they will be supposed the
same whenever the corresponding variables are made equal. This equality will be
denoted by writing, for instance,
\[
\square V V' V''\ldots
\]
to represent the value assumed by
\[
\square V_1V_2V_3V_4\ldots
\]
when after the differentiations
\[
x_1, y_1 = x_2, y_2 = x_3, y_3 = x, y,
\]
\[
x_4, y_4 = x', y';
\]
\&c.

It is easy to determine the general form of $\square U$. To do this, writing for shortness
\[
\alpha + \beta + \gamma + \ldots = f_1,
\]
\[
\alpha + \beta' + \gamma' + \ldots = f_2,
\]
\[
\beta + \beta' + \gamma'' + \ldots = f_3,
\]
\&c.,
\[
N = (-)^{a+a'+a''+\ldots+l+l'+l''+\ldots}
\frac{(\alpha)!}{(\overline{r})!}
\frac{(\beta)!}{(s)!}
\frac{(\gamma)!}{(t)!}\ldots
\frac{(\delta)!}{(\overline{r}')!}
\frac{(\gamma')!}{(s')!}\ldots
\frac{(\beta'')!}{(\overline{r}'')!}\ldots
\]
\[
k^{r+s+t+\ldots}\,V = \delta_x^r\delta_y^{s+\ldots}\,V \quad\text{or}\quad V^{r,s},
\]
the general term is
\[
N \cdot V_1^{f_1}\,V_2^{f_2}\,\ldots\,V_p^{f_p},
\]
where $r$, $s$, $t$, \ldots $r'$, $s'$, $r''$, \ldots extend from 0 to $s$, $t$, \ldots $\beta'$, $\gamma'$, \ldots\ respectively. It
would be easy to change this general term in a way similar to that which will be
employed presently, for the particular case of $\square V_1V_2V_3$.

If several of the functions become identical, that is when some of the letters
$f$ are equivalent, it is clear that the derivative $\square U$ refers to a certain number of
functions $V_1$, $V_2$,\ldots the same as before, of different, the variables $x$, $y$; $x'$, $y'$;\ldots\ respectively. It
that this derivative is homogeneous in the degrees $\theta_1$, $\theta_1'$,\ldots\ with respect to the
differential coefficients of the orders $f_1$, $f_1'$,\ldots\ (consequently homogeneous in the
the order $\theta_1+\theta_1'+\ldots=p$ suppose. In general, only a single function will be considered,

\newpage

and it will be assumed that $\square U$ only contains the differential coefficients of the
$f^{\text{th}}$ order. In this case, the derivative is said to be of the degree $p$ and of the
order $f$. The most convenient classification is by degrees, rather than by orders.

Commencing with the simplest case, that of functions of the second order (and
writing $V$, $W$ instead of $V_1$, $V_2$), we have
\[
\square VW = \overline{12}^a\,V'W';
\]
(where $\xi_1$, $\eta_1$ apply to $V$ and $\xi_2$, $\eta_2$ to $W$). This will be constantly represented in
the sequel by the notation
\[
\overline{12}^a\,VW = B_a\,(V,\ W).
\]
Hence, writing
\[
\delta_x V = V_r,\quad \delta_{a-r}V = V^{a-r}\ldots,
\]
we have
\[
B_a\,(V,\ W) = V_a\cdot W_a - \left[\frac{a}{1}\right]V_a^{a-1}\cdot W_{a-1} + \ldots;
\]
and in particular, according as $a$ is odd or even,
\[
B_a\,(V,\ V) = 0,
\]
\[
\tfrac{1}{2}B_a\,(V,\ V) = V_a^a - \frac{a}{1}V_a^{a-1}\cdot V_{a-1} + \ldots,
\]
continued to the term which contains $V_a^{a/2}V_{a-a/2}$, the coefficient of this last term
being divided by two.

Thus, for the functions $\tfrac{1}{2}(ax^2 + 2bxy + cy^2)$, $\tfrac{1}{2}(ax^4 + 4bx^3y + 6cx^2y^2 + 4dxy^3 + ey^4)$, \&c.,
if $a$ be made equal to $2$, $4$, \&c.\ respectively, we have the constant derivatives
\[
ac - b^2,
\]
\[
ae - 4bd + 3c^2,
\]
\[
ag - 6bf + 15ce - 10d^2,
\]
\[
ai - 8bh + 28cg - 56df + 35e^2,
\]
\[
\vdots
\]
which have all of them the property of remaining unaltered, \textit{\`a un facteur pr\`es}, when
the variables are transformed by means of $x = \lambda\bar{x} + \mu\bar{y}$, $y = \lambda'\bar{x} + \mu'\bar{y}$. Thus, for instance,
if these equations give
\[
ac - b^2 = (\lambda\mu' - \lambda'\mu)^2\,(ac - b^2);
\]
then
\[
ac - b^2 = (\lambda\mu' - \lambda'\mu)^2\,(ac - b^2);
\]
and so on. This is the general theory, which we call to mind for the case of these
constant derivatives.

\newpage

The above functions may be transformed by means of the identical equation
\[
B_n\,(V,\ W) = \overline{12}^{n-2}\,B_2\,(V,\ W),
\]
to make use of which, it is only necessary to remark the general formula
\[
\xi_1^r\eta_1^s\eta_2^r B_2\,(V,\ W) = B_2\,(\xi^r\eta^s V,\ \xi^r\eta^s W).
\]
Thus, if $k=1$, we obtain for the above series, the new forms
\begin{align*}
&ac - b^2, \\
&\tfrac{1}{2}\{(ae - bd) - 3(bd - c^2)\}, \\
&(ag - bf) - 5(bf - ce) + 10(ce - d^2), \\
&\tfrac{1}{2}\{(ai - bh) - 7(bh - cg) + 21(cg - df) - 35(df - e^2)\},
\end{align*}
the law of which is evident. This shows also that these functions may be linearly
expressed by means of the series of determinants
\[
\left|
\begin{array}{cc}
a,&b \\
b,&c
\end{array}
\right|,\quad
\left|
\begin{array}{ccc}
a,&b,&c \\
b,&c,&d
\end{array}
\right|,\quad\&\text{c.}
\]

We may also immediately deduce from them the derivatives $B$ which relate to two
functions. For example, for functions of the sixth order this is
\[
ag' + a'g - 6(bf' + b'f) + 15(ce' + c'e) - 20\,dd',
\]
which has an obvious connection with
\[
ag - 6bf + 15ce - 10d^2;
\]
and the same is the case for functions of any order.

The following theorem is easily verified; but I am unacquainted with the general
theory to which it belongs.

``If $U$, $V$ are any functions of the second order, and $W = \lambda U + \mu V$; then
\[
B_2'\,[B_2\,(W,\ W),\ B_1\,(W,\ W)] = 0,
\]
(where $B_2'$ relates to $\lambda$, $\mu$) is the same that would be obtained by the elimination of
$x$, $y$ between $U=0$, $V=0$.''\footnote{Not given with the present paper.}

In fact this becomes
\[
4\,(ac - b^2)(a'c' - b'^2) - (ac' + a'c - 2bb')^2 = 0,
\]
which is one of the forms under which the result of the elimination of the variables
from two quadratic equations may be written. This is a result for which I am
indebted to Mr~Boole.

\end{document}
